\section{Related Work}
\label{sec:related}

\subsection{Foundation Audio Models and Cross-Cultural MIR}

Foundation audio models provide strong upstream representations for music information retrieval, but they do not by themselves specify how cross-cultural recommendation should be structured. MERT established that large-scale self-supervised training can yield strong generic music embeddings for downstream MIR tasks \cite{li2023mert}. CultureMERT extends this line through continual pre-training on a multicultural mixture and reports improved performance on non-Western auto-tagging settings \cite{kanatas2025culturemert}. We are directly inspired by this literature in one sense: both CultureMERT and Gemini are treated here as frozen embedding backbones rather than handcrafted features. Unlike these works, however, our contribution does not lie in designing or adapting the upstream encoder itself. The question we ask is narrower and downstream-facing: once strong audio embeddings already exist, what additional structure is needed for culturally calibrated recommendation?

Recent cross-cultural MIR benchmarks sharpen this distinction. GlobalMood argues for culturally grounded evaluation resources rather than assuming that a single label ontology transfers universally \cite{lee2025globalmood}. Likewise, recent evaluation of foundation models on world-music corpora shows that strong audio encoders can still degrade on underrepresented repertoires \cite{papaioannou2025universal}. These studies motivate our problem setting and provide evidence that cross-cultural generalization remains difficult. Unlike them, however, we do not focus on benchmark construction alone or on the absolute strength of a new representation model. Instead, we build a backbone-agnostic downstream recommendation stack that explicitly addresses target-culture matching, minority exposure, and human correction on top of existing backbones.

\subsection{Factorized Representation, Invariance, and Calibration-Aware Recommendation}

Another line of work studies whether useful representations can be factorized into more interpretable or more robust components. \(\beta\)-VAE introduces a constrained variational objective that encourages disentangled latent factors \cite{higgins2017betavae}, while FactorVAE improves factor separation through an explicit total-correlation penalty \cite{kim2018factorvae}. Our method is inspired by this factorized-latent lineage, but our goal differs in two important ways. First, we do not claim perfect semantic disentanglement. Second, the factorization is not an end in itself: it is used to separate culturally specific style from the affective-functional subspace that supports cross-cultural recommendation. In that sense, our method is closer to task-useful factorization than to generic unsupervised representation discovery.

Invariant learning provides a second source of inspiration. Domain-adversarial training encourages a representation to remain informative for the main task while becoming less discriminative for nuisance domains \cite{ganin2016dann}. We adopt this idea to reduce direct culture leakage in the latent subspace used for matching, but we also differ from classical invariance-based formulations. Rather than removing all culture-specific information, we intentionally preserve stylistic-cultural variation in \(z_s\) and shape \(z_a\) as the shared matching space. This selective use of invariance is important in cross-cultural recommendation, where cultural style is not noise to be erased but structure to be routed appropriately.

Our reranking layer is also related to calibrated recommendation and efficient transport-based matching. Calibrated recommendation argues that the output distribution of a recommender should better reflect meaningful user interests rather than optimize relevance alone \cite{steck2018calibrated}. Sinkhorn regularization makes optimal transport practical for distributional alignment \cite{cuturi2013sinkhorn}. We are inspired by both ideas, but our use case is different from either one in isolation. Unlike generic calibration methods, our calibration-aware reranking jointly controls target-culture affinity, minority exposure at k, source diversity, and novelty in a cross-cultural setting. Unlike standalone OT methods, transport here is only one component in a broader downstream stack that includes factorized representation learning and late-stage human feedback.

\subsection{Serendipity, Diversity, and Human-Centered MIR Applications}

Cross-cultural recommendation also connects to work that treats discovery, diversity, and human input as first-class design goals. Auralist is a classic example of serendipity-oriented music recommendation, emphasizing that recommendation quality should combine relevance with unexpectedness rather than optimize accuracy alone \cite{zhang2012auralist}. Diversity-by-design work makes a related argument at the recommender-system level, namely that exposure and catalog balance should be engineered into the system rather than left as accidental by-products \cite{porcaro2021diversity}. We are clearly inspired by this perspective. At the same time, our setting differs from both strands because we do not optimize catalog-level novelty alone. The relevant problem is calibrated cross-cultural transfer, where minority exposure at k and target-culture affinity must be balanced jointly instead of being collapsed into a single diversity score.

The same difference appears in our feedback design. Active learning classically asks which examples should be labeled next under a limited annotation budget \cite{settles2009active}. Recent human-centered MIR work on underrepresented traditions further shows that annotation interfaces and expert knowledge cannot be treated as interchangeable commodities \cite{pinto2025maracatu}. Responsible MIR scholarship makes the broader point that cultural coverage, annotation authority, and system goals must be aligned if MIR is to avoid reproducing structural asymmetries \cite{holzapfel2018ethical,huang2023beyonddatasets}. Our PAL design is inspired by these ideas, but it also differs from them in scope and mechanism. Unlike generic active learning, we do not target large-scale class labeling; instead, we rank uncertain cross-cultural boundary cases and convert expert judgments into pairwise constraints for stage-3 retraining. Unlike responsible MIR position papers, we do not stop at normative critique; we instantiate a runnable PAL-ready workflow that can be attached to a recommendation system.

Across these lines of work, the common limitation is that upstream representation learning, latent regularization, calibration, and human feedback are usually treated as separate problems. Existing work either strengthens music embeddings, improves latent objectives, or argues for diversity- and human-centered MIR evaluation, but it still leaves one gap open: how to build a backbone-agnostic downstream recommendation stack that jointly factorizes cross-cultural similarity, exposes controllable calibration trade-offs, and incorporates targeted human correction under source-confound-aware evaluation. That gap is the entry point of this paper.
